\section*{Основные определения, термины и сокращения}
\addcontentsline{toc}{section}{Основные определения, термины и сокращения}

\begin{itemize}
\item \textbf{Поездка} --- сущность планирования, объединяющая даты, участников, идеи, маршрут и расходы.
\item \textbf{Участник} --- пользователь, присоединенный к поездке и участвующий в совместных действиях.
\item \textbf{Владелец} --- пользователь, который создал поездку и управляет доступом и модерацией.
\item \textbf{Идея} --- пользовательское предложение активности или места до включения в маршрут.
\item \textbf{Активность} --- элемент маршрута, назначенный на конкретный день поездки.
\item \textbf{Расход} --- запись о планируемой или уже произведенной трате в рамках поездки.
\item \textbf{Баланс} --- агрегированное представление финансовых обязательств участников поездки.
\item \textbf{Офлайн-синхронизация} --- механизм локального сохранения действий при отсутствии сети с последующей отправкой изменений на сервер.
\item \textbf{AI-предложение} --- автоматически сгенерированная рекомендация, которую пользователь может сохранить в список идей.
\end{itemize}
